% ============================================================
%  Paper A — Main Text
%  "Plant Science is Reorganizing Around Climate-Relevant
%   Crops: Evidence from 4.9 Million Publications"
% ============================================================

% ── Introduction ─────────────────────────────────────────
\section*{Introduction}

Plant science underpins the food, fibre, and fuel supply of eight billion
people, yet it now faces twin pressures. From without, climate change is
altering growing seasons, amplifying drought and heat stress, and shifting
pathogen distributions across every agricultural system~\cite{ipcc2022}. From
within, molecular and computational tools --- high-throughput sequencing, gene
editing, and machine learning --- have for the first time made it feasible to
engineer complex agronomic traits directly in crops rather than first in model
organisms~\cite{jackson2011}. Whether the field has actually redirected its
attention toward this challenge is an empirical question that bibliometrics is
positioned to answer.

Twentieth-century plant biology was organized around model organisms, most
conspicuously \textit{Arabidopsis thaliana}, whose compact genome, rapid
generation time, and rich mutant collections made it an ideal substrate for
dissecting development, hormone signalling, and stress
response~\cite{somerville1986,meyerowitz1989}. Yet translation from
\textit{Arabidopsis} to wheat or soybean routinely required years of extra work,
and many agronomic traits --- grain quality, nitrogen-use efficiency, drought
tolerance in deep-rooted species --- had no tractable \textit{Arabidopsis}
equivalent~\cite{feuillet2011,voss2015}. CRISPR-Cas9 and pan-genome resources
lowered the barrier to working directly in crops, raising the hypothesis that the
field was shifting from model-organism-centred discovery toward crop-centred
translational research.

We test this by analysing \num{4.9} million plant science publications
(1900--2024; the modern window 1990--2024 is used for field-growth and
organism-share analyses). Across six independent lines of evidence --- growth has
plateaued, \textit{Arabidopsis} share has declined, sleeping beauties on minor
crops are awakening, CRISPR diffusion is unusually rapid, crop-genomics concept
marriages are proliferating, and orphan crops remain under-served --- plant
science emerges as being in a genuine paradigm shift, but one that must accelerate
to meet the demands of a warming world. A temporal placebo analysis confirms the
signal is specific to recent decades (permutation $p = 0.001$), ruling out
background publication growth.

% ── Results ──────────────────────────────────────────────
\section*{Results}

\subsection*{Field growth has reached a logistic plateau}

Annual publication output in the deduplicated plant-science corpus grew
from approximately 40,000 papers per year in 1990 to a peak of roughly
199,000 papers per year in 2019, with 2022--2024 counts at or slightly
below that peak.
When we fit both exponential and logistic growth models to this trajectory
(works\_clean, 1990--2024), the logistic model provides a significantly
better description of the data ($\Delta$AIC $= 31.2$, Figure~\ref{fig:FA1}).
We also compared logistic against Gompertz and Richards (generalised
logistic) sigmoidal alternatives; Richards fits better
($\Delta\mathrm{AICc} = 11.5$ relative to logistic) and Gompertz worse
($\Delta\mathrm{AICc} = 18.9$), with all three models agreeing that the
field has reached the plateau phase (Supplementary Note~\ref{sec:growth-model-aic}).
We retain the logistic parameterisation in the main text for
interpretability of $K$ and $t_{\mathrm{mid}}$.
The fitted logistic carrying capacity is $K = 217{,}781$ papers per year,
and the inflection point --- the year at which year-on-year growth was
fastest --- falls in 2006.
Since 2006, year-on-year growth rates have declined monotonically toward
zero, and the most recent five-year window (2020--2024) shows near-zero
net growth after correcting for indexing lag.

This plateau admits three interpretations: genuine saturation of research effort;
a compositional shift in which new topics enter as fast as older ones decline,
holding aggregate counts stable while the internal structure changes; or a
database-coverage ceiling. The analyses that follow point to compositional change
rather than ceiling effects or saturation, and motivate what follows.

\begin{figure}[H]
  \centering
  \includegraphics[width=0.85\textwidth]{FA1_growth_logistic}
  \caption{\textbf{Logistic growth plateau in plant science publications.}
    Annual publication counts (\texttt{works\_clean}, 1990--2024,
    $n \approx 3{,}983{,}000$ deduplicated) shown as bars, overlaid
    with the fitted logistic growth model (solid black curve). The
    logistic fit yields carrying capacity $K = 217{,}781$ papers/year
    with inflection point in 2006 ($\Delta\mathrm{AIC} = 31.2$ vs an
    exponential baseline; Gompertz and Richards alternatives
    discussed in Supplementary Note~\ref{sec:growth-model-aic}).
    Horizontal dotted line marks the fitted carrying capacity
    $K$; vertical dashed line marks the 2006 logistic inflection.
    Bars are colour-coded by era: pre-1990 pre-genomics (grey),
    1990--2006 acceleration (teal), 2006--2024 post-inflection
    approach to $K$ (green). The 2019 observed peak
    ($\sim 199{,}000$ papers/year) is approximately $91\%$ of $K$.}
  \label{fig:FA1}
\end{figure}

\subsection*{Model organism decline and crop species ascent}

Within the plateauing total, the relative share of publications mentioning
\textit{Arabidopsis thaliana} rose from 4.1\% in 1990 to a maximum of
11.3\% in 2012 before declining to 8.7\% in 2024 (Extended Data Fig.~\ref{fig:FA2}).
This trajectory mirrors the logistic field-level plateau almost exactly. The
same year that aggregate growth slowed marks the zenith of
\textit{Arabidopsis} dominance.
Over the same period, the combined share of rice (\textit{Oryza sativa}),
wheat (\textit{Triticum aestivum}), and maize (\textit{Zea mays}) rose from
21.4\% in 1990 to 34.1\% in 2024.

The transition is not merely a reallocation among established species.
Soybean, sorghum, and cassava --- crops of particular importance for protein
supply and food security in low-income countries --- all show accelerating
growth rates after 2015.
Notably, the acceleration in these crops coincides with the rapid adoption of
CRISPR-Cas9 (see below), suggesting that the availability of a crop-agnostic
editing platform removed a key barrier to working in genetically recalcitrant
species.
Hierarchical clustering of organism-level publication trajectories identifies
three distinct groups: a declining model-organism cluster led by
\textit{Arabidopsis}; a stable staple-crop cluster (rice, wheat, maize); and
an emerging minor-crop cluster showing exponential growth from a low base
(per-organism growth rates in Supplementary Fig.~\ref{fig:S1}).

The share decline is striking because \textit{Arabidopsis} absolute output kept
growing through approximately 2018: other organisms are simply growing faster,
and what is changing is \textit{Arabidopsis}'s relative centrality, not its
productivity.

% [Fig FA2 moved to Extended Data]

\subsection*{\textit{Arabidopsis} research is consolidating as crop science's foundational methodological substrate}

The organism-share decline documented above invites a pessimistic
reading, that \textit{Arabidopsis} is becoming irrelevant to a
crop-centred field. A direct analysis of citation flow rejects that
reading.

We define the instrumentalization ratio
$r_{\mathrm{inst}}(Y) = [\text{C}{\to}\text{A}\text{ share in year }Y] / [\text{A}{\to}\text{A}\text{ share in year }Y]$,
where C is the rice/wheat/maize crop triad and A is \textit{Arabidopsis}
(\textit{Methods}; paper\_a\_organism.csv confidence $\geq 0.7$).
Between 1995 and 2016, $r_{\mathrm{inst}}$ rose from 0.018 to 0.161
--- a nine-fold monotonic increase --- before plateauing and modestly
declining to 0.131 by 2024 (Figure~\ref{fig:M9}a). In parallel, the
median lag from publication of a top-100 most-cited
\textit{Arabidopsis} paper to its first citation by a rice, wheat, or
maize paper collapsed from 6.5 years (1995 cohort) to a stable
two-year floor by the 2002 cohort and onward (Figure~\ref{fig:M9}b).
Confidence intervals are 95\% bootstrap percentile (1{,}000 resamples)
throughout.

Taken together, the two curves describe \textit{Arabidopsis}
transitioning from a standalone model-organism paradigm to a
foundational methodological substrate --- a lookup library --- from
which crop science launches translational work. This reframes the
organism-share decline as consolidation, not retreat. The field has
not abandoned \textit{Arabidopsis}, it has completed the work of
\textit{Arabidopsis} as a topic and is now using it as a tool. The
2010--14 peak in $r_{\mathrm{inst}}$ coincides with an independent
2010--14 peak in coauthor-network paradigm alignment reported in
Supplementary Note~\ref{sec:coauthor-alignment}, providing convergent evidence that
the consolidation reached steady state in the mid-2010s.
Supplementary analyses further show that
\textit{Arabidopsis}-specialist authors retain a structural citation
advantage (Supplementary Note~\ref{sec:bridge-scientists}), consistent with their
work functioning as the canonical reference library for downstream
crop research.

\begin{figure}[H]
  \centering
  \begin{subfigure}[b]{0.48\textwidth}
    \includegraphics[width=\textwidth]{fig_instrumentalization}
    \caption{}
    \label{fig:M9a}
  \end{subfigure}
  \hfill
  \begin{subfigure}[b]{0.48\textwidth}
    \includegraphics[width=\textwidth]{fig_reverse_translation_lag}
    \caption{}
    \label{fig:M9b}
  \end{subfigure}
  \caption{\textbf{\textit{Arabidopsis} consolidates as crop science's
    methodological substrate.}
    (a) Instrumentalization ratio
    $r_{\mathrm{inst}} = [\text{C}{\to}\text{A}] / [\text{A}{\to}\text{A}]$
    between 1995 and 2024. Crop$\to$\textit{Arabidopsis} citations grew
    from 1.8\% of \textit{Arabidopsis} self-citation in 1995 to 16.1\%
    in 2016 before plateauing.
    (b) Reverse-translation lag: median years to first rice/wheat/maize
    citation of a top-100 \textit{Arabidopsis} paper, by cohort year.
    Lag collapsed from 6.5 years (1995 cohort) to a 2-year floor
    (2002 cohort and onward).}
  \label{fig:M9}
\end{figure}

\subsection*{Climate shocks drive reactive publication at 6- and 18-month lags}

Does the reorganization track the climate pressure that motivates it? We ran a
pooled difference-in-differences event study on 15 pre-registered climate events
(EM-DAT, NOAA Billion-Dollar Disasters, and IPCC/UNFCCC policy events;
\textit{Online Methods}, Supplementary Note~\ref{sec:climate-shocks}), comparing
detrended monthly publication counts on affected taxa or policy keywords against
matched controls over $[t-36, +36]$ months. The pooled residual shows two
significant response windows: a fast spike at lag $+6$ months (commentary and
reviews) and a main spike at $+18$ months (empirical studies), the latter
$+376.8$ papers per month above control (95\% CI $[+28.6, +773.8]$). Against a
field publishing roughly $15{,}000$ papers per month, this $\sim$2--3\%
redirection of output falls in exactly the months that translational-research
funding cycles would predict.

\subsection*{Corpus-wide organism classification and paradigm analysis}

To complement keyword-based assignment, we applied SPECTER2 embeddings and
zero-shot classification to the full classified corpus of \num{2.77} million
papers (\textit{Online Methods}). Among organism-specific papers (66\% of the
corpus carries no organism signal), rice (3.9\%), wheat (3.9\%), and maize
(3.6\%) collectively exceed \textit{Arabidopsis} (2.3\%) in recent share
(Figure~\ref{fig:M1}) --- a direct, model-free crossover confirming the
structural transition (head-to-head in Extended Data Fig.~\ref{fig:M2}).

\begin{figure}[H]
  \centering
  \includegraphics[width=0.85\textwidth]{M1_organism_share_timeseries}
  \caption{\textbf{Organism share timeseries from zero-shot classification
    (2.77 million papers).}
    Relative publication share for each organism category 1990--2024, derived
    from corpus-wide zero-shot classification.
    Non-specific papers (66\% of corpus) are excluded for clarity.
    Note that rice, wheat, and maize collectively surpass \textit{Arabidopsis}
    share in recent years.}
  \label{fig:M1}
\end{figure}

% [Fig M2 moved to Extended Data]

A chi-square test of paradigm (applied vs.\ fundamental) against organism is
highly heterogeneous ($\chi^2 = 12{,}361$, $df = 9$, $p \approx 0$): 80\% of the
corpus is applied, but the fraction varies --- \textit{Arabidopsis} skews applied
at 95.6\%, while maize (78.2\%) and other model organisms (77.9\%) retain more
fundamental research (Extended Data Fig.~\ref{fig:M3}; citation impact by
paradigm in Supplementary Fig.~\ref{fig:S2}).

% [Fig M3 moved to Extended Data]

\subsection*{Topic landscape --- 30 macro-themes, convergence, and emerging research}

BERTopic applied to SPECTER2 embeddings of the full classified corpus yields
1,632 fine-grained topics, which we aggregate into 30 coherent macro-themes
by hierarchical merging (Extended Data Fig.~\ref{fig:M4}).
Across the 2020s, 9 macro-themes are classified as emerging (increasing
publication share) and 7 as declining; 14 are stable.
Which organisms concentrate in which macro-themes is quantified by the
organism--topic enrichment matrix (Extended Data Fig.~\ref{fig:M6}).
Shannon entropy of the macro-theme distribution has declined from a mean of
1.82 bits in the 1990s to 1.14 bits in the 2020s ($\Delta = -0.60$ bits),
a 37\% reduction indicating substantial concentration of research attention
around a narrower set of themes (Extended Data Fig.~\ref{fig:M5}; annual
entropy series in Supplementary Fig.~\ref{fig:S3}).

% [Fig M4 moved to Extended Data]

% [Fig M5 moved to Extended Data]

% [Fig M6 moved to Extended Data]

\subsection*{Semantic drift --- crop species converging toward Arabidopsis}

Using SPECTER2 embeddings, we track each organism's conceptual profile as the
cosine distance of its centroid from \textit{Arabidopsis} in rolling five-year
windows (Extended Data Fig.~\ref{fig:M7}). Cotton and rice converge most strongly
toward \textit{Arabidopsis} concept space over 2010--2024, consistent with
adoption of tools originally developed in the model organism, whereas
non-specific papers remain most distant ($d = 0.047$, 2020--2024). This semantic
convergence complements the citation-flow findings: crop papers are not merely
citing \textit{Arabidopsis} but increasingly resemble it in concepts and methods.

% [Fig M7 moved to Extended Data]

\subsection*{Disruption index and citation dynamics by organism}

We compute the CD disruption index~\cite{funk2017disruption} at decade resolution
for each organism (Figure~\ref{fig:S6}). \textit{Arabidopsis} has the most
negative mean value (CD\,$= -0.001$), consistent with a consolidating literature,
whereas crop and non-specific papers are more disruptive (maize $0.095$, soybean
$0.102$, cotton $0.099$, non-specific $0.107$). These organism-level differences
are a null result at this resolution: the Kruskal--Wallis test is not significant
($H = 12.9$, $p = 0.17$, $n_{\mathrm{groups}} = 10$), and although the direction is
consistent across all 10 groups, the standardised effect is negligible
($\eta^2 = 0.0001$) despite $99.6\%$ power ($0.996$) on $N = 30{,}462$
organism-decade observations. We therefore exclude the CD index from the main
evidence chain and report it for completeness. \textit{Arabidopsis} also shows the
highest self-citation rate (61.6\%; cross-citation matrix in Supplementary
Fig.~\ref{fig:S9}) and international collaboration rate (18.9\%; Supplementary
Fig.~\ref{fig:S10}), while crop organisms have lower, more variable rates
(rice 13.8\%, wheat 15.2\%).

\subsection*{CRISPR adoption lags differ by organism}

CRISPR-Cas9 diffusion is uneven across organisms (Extended Data
Fig.~\ref{fig:M8}): \textit{Arabidopsis}, rice, and wheat adopt earliest and
fastest, maize lags by one to two years (consistent with harder transformation),
and minor crops and non-specific papers adopt slowest --- CRISPR
democratisation, while real, has not yet reached the most neglected species.

% [Fig M8 moved to Extended Data]

\subsection*{Sleeping beauties on minor crops are awakening}

A sleeping beauty is a paper that receives negligible citations for an extended
dormancy before a sudden, sustained burst~\cite{ke2015sleeping}. Applying the Ke
et al.\ index $B$ to the full corpus, we identify \num{40} papers on minor or
under-studied crops meeting strict criteria (dormancy $\geq 8$ years, awakening
amplitude $\geq 10\times$), spanning 23 crop species including teff, quinoa,
amaranth, cowpea, and breadfruit --- long peripheral to mainstream plant biology.
Their awakenings are strikingly non-uniform: 35 of the 40 (87.5\%) fall in
2019--2022 (Figure~\ref{fig:FA3}), with median dormancy 11.4 years and median
awakening magnitude $17.3\times$ (highest-index papers in Supplementary
Table~\ref{tab:S2}). The citing papers are empirical research programmes rather
than reviews, treating the dormant work as a foundational reference once concern
about diet diversification, climate adaptation, and food sovereignty made it
relevant. With $n = 40$ from a 4.9-million corpus this is a case study of
specific minor-crop taxa, not field-wide evidence; the narrow statistical claim
is that these papers' awakening times cluster --- the probability of $\geq 35$
awakenings in a four-year window, given the historical rate across the 34-year
corpus, is $p < 0.001$ (permutation test, $10^5$ iterations). Major crops (rice,
wheat, maize) show no comparable clustering.

\begin{figure}[H]
  \centering
  \includegraphics[width=0.85\textwidth]{FA3_sleeping_beauties}
  \caption{\textbf{Sleeping beauty awakenings cluster in 2019--2022.}
    Histogram of awakening year for the 40 sleeping-beauty papers on minor
    crops identified in the corpus.
    Inset table lists the five papers with the highest sleeping-beauty index
    $B$, including crop species, dormancy length, and awakening magnitude.
    The 2019--2022 cluster (shaded) contains 87.5\% of all awakenings.}
  \label{fig:FA3}
\end{figure}

\subsection*{Method diffusion is accelerating the paradigm shift}

Methodological innovation is historically a leading indicator of paradigm shifts
in biology~\cite{kuhn1962,hegarty2008}. Fitting sigmoid adoption models to five
method categories --- genetic mapping, next-generation sequencing, RNA
interference, CRISPR-Cas9, and machine learning (Figure~\ref{fig:FA4}) ---
CRISPR-Cas9 shows the fastest diffusion in the dataset: 9.2 years from its first
plant-science application (2013) to crossing the 5\% publication-share threshold,
versus 16.8 years for NGS, 14.1 for RNAi, and 22.3 for QTL mapping. Its
accessibility (no specialised equipment) and crop-agnostic applicability --- over
60 plant species within a decade --- underlie this speed. Machine learning does
not fit a classical sigmoid but grows super-exponentially, a 39-fold rise since
2010; the plant-science-specific rate ($r = 0.31$/year) exceeds the cross-field
baseline ($r = 0.19$/year), with phenotyping accounting for most of it. Together,
CRISPR and ML remove the two historical barriers to crop-centred research ---
editing complex polyploid genomes and measuring quantitative traits at scale ---
so their concurrent adoption is mechanistically consistent with the
organism-level shifts above.

\begin{figure}[H]
  \centering
  \includegraphics[width=0.85\textwidth]{FA4_method_diffusion}
  \caption{\textbf{Method diffusion S-curves in plant science.}
    Publication share (percentage of annual plant science output) for five
    methodological categories, 1990--2024, with fitted logistic adoption
    curves (solid lines).
    Numbers in brackets indicate time to reach 5\% share from year of
    introduction.
    Machine learning (gold) is shown on a separate secondary axis due to
    its super-exponential growth pattern (inset, log scale).}
  \label{fig:FA4}
\end{figure}

\subsection*{Concept marriages define a new intellectual epoch}

Intellectual breakthroughs often occur at the intersection of previously
separate research traditions --- what we term ``concept marriages.'' We detect
these by tracking Jaccard co-occurrence between keyword sets across five-year
windows, flagging combinations that cross a threshold above the background rate
(\textit{Online Methods}); this lexical proxy indexes semantic clusters in the
keyword network rather than adjudicated paradigms. Hierarchical clustering of
the co-occurrence network partitions the corpus into six intellectual epochs
(Extended Data Fig.~\ref{fig:FA5}), culminating in a current epoch defined by
the marriage of crop genomics with precision breeding and machine-learning
phenotyping --- a combination impractical a decade earlier. This epoch is the
most crop-centred (wheat, rice, maize, soybean; yield, drought tolerance,
nitrogen use) and shows the highest concept-marriage formation rate in the
dataset (3.7 new marriages per year vs.\ a historical mean of 1.9). The same
reorganisation appears in the field's citation canon, whose top-100 reference
list turned over substantially over the period
(Supplementary Note~\ref{sec:canon-churn}).

% [Fig FA5 moved to Extended Data]

\subsection*{The orphan organism gap persists}

Despite these signals of transition, a systematic gap persists between crop
species' economic importance and the research attention they receive. We
quantify it as the studentised residual from a log-log regression of publication
count on a composite FAO economic-importance index (production value, area
harvested, and number of countries where the crop is a staple;
\textit{Online Methods}); this denominator is normative --- alternatives such as
phylogenetic diversity or food-security criticality would reweight the ranking.
Oil palm shows the most extreme negative residual ($z = -5.21$), followed by
sugarcane ($z = -4.03$), cassava ($z = -3.87$), and yam ($z = -3.12$)
(Figure~\ref{fig:FA6}; Supplementary Table~\ref{tab:S1}), whereas
\textit{Arabidopsis} ($z = +8.94$) and tobacco ($z = +4.71$) sit far above the
line. Rice, wheat, and maize lie close to it, so the staple-crop shift is
redistributing attention toward these species but not yet reaching the orphan
crops with the greatest unmet need. This gap is not new~\cite{nair2010}; what our
data add is temporal resolution --- the oil-palm and sugarcane gaps were already
large in 1990 and have widened through 2024. Closing this two-speed divergence
will require targeted funding and career incentives that the market forces
reshaping the top of the distribution do not provide.

\begin{figure}[H]
  \centering
  \includegraphics[width=0.85\textwidth]{FA6_orphan_organisms}
  \caption{\textbf{The orphan organism research-attention gap.}
    Log-log scatter of publication count (2020--2024) against composite
    economic importance index for 120 economically significant plant species.
    Dashed line: ordinary least squares regression ($R^2 = 0.61$).
    Points coloured by attention-gap $z$-score (blue: under-studied;
    red: over-studied relative to economic importance).
    Selected species are labelled.
    Inset: oil palm gap trajectory 1990--2024, showing persistent negative
    residuals.}
  \label{fig:FA6}
\end{figure}

% ── Discussion ───────────────────────────────────────────
\section*{Discussion}

Nine independent lines of evidence --- logistic growth saturation, model-organism
decline, sleeping-beauty awakening clusters, accelerated method diffusion, a new
crop-centred concept epoch, a persistent orphan-organism gap, corpus-wide
classification confirming crop species now collectively exceed \textit{Arabidopsis}
in share, a topical-entropy decline (1.82\,$\to$\,1.14 bits), and semantic
convergence of crop literature toward \textit{Arabidopsis} --- converge on one
conclusion: plant science is undergoing a structural reorganisation from
model-organism-centred discovery toward crop-centred, climate-relevant research.
Their methodological diversity (growth modelling, citation dynamics, diffusion
fitting, semantic-network analysis, and attention-gap regression) makes any single
artefact unlikely, and a temporal placebo test confirms the signal is specific to
the recent era (Supplementary Note~\ref{sec:temporal-placebo}).

The year 2012 is a qualitative nexus: the peak of \textit{Arabidopsis} relative
share, the arrival of CRISPR-Cas9 in plant biology, and the deep-learning
revolution in image analysis. Field growth had already begun slowing at the 2006
logistic inflection (Fig.~\ref{fig:FA1}), but only around 2012 did the agenda
redirect \emph{in kind} rather than merely saturate in volume. Two drivers appear
to act in concert: the CRISPR-Cas9 platform, whose 9.2-year diffusion is the
fastest we measure and whose crop-agnostic applicability suited it to the
transition; and climate change, whose intensifying disruption through the
2010s--2020s and the 2019--2022 sleeping-beauty cluster are consistent with
crisis-driven recognition of neglected science. Our corpus-level metrics cannot
adjudicate this causal reading against alternatives such as demographic cohort
effects or funding-agency re-prioritisation~\cite{kuhn1962,price1963}.

Consistent with model-organism-to-applied transitions documented in
neuroscience~\cite{bhatt2020}, cancer biology~\cite{rhodes2021}, and
microbiology~\cite{sunagawa2015}, the shift is not a clean replacement:
\textit{Arabidopsis} absolute output kept growing through 2018 even as its share
fell, and current-epoch concept marriages combine crop genomics with model-system
insight. The transition is better read as expansion and rebalancing --- the
consolidation of \textit{Arabidopsis} into crop science's methodological substrate
--- than as decline.

Several limitations apply. Publication counts are a noisy proxy for effort, and
databases under-represent non-English and applied-agricultural literature; the
keyword- and abstract-based analyses may miss full-text content; the FAO-based
attention-gap denominator is contestable; and organism-level disruption (CD)
differences are not significant at the decade level ($p = 0.17$) and should be
read as directional. Topic analyses cover the 46.1\% of classified papers with a
valid BERTopic assignment, though the excluded outliers do not differ appreciably
in organism or decade composition (\textit{Online Methods}). Robustness details
for the climate-event study, the paradigm-bridging-author analysis, and the
coauthor-network analysis are given in Supplementary
Notes~\ref{sec:climate-shocks},~\ref{sec:bridge-scientists},
and~\ref{sec:coauthor-alignment}. Our data end in 2024; whether the crop-centred
epoch consolidates, accelerates, or stalls only future data can resolve. What the
present evidence establishes clearly is that the transition is real, measurable
across multiple independent dimensions, and that closing the orphan-organism gap
will require deliberate policy intervention rather than a continuation of current
trends.

\subsection*{Policy implications}

Three implications follow. First, funding calls organised under ``model
organism'' headers risk under-representing the crop-centred research that now
dominates the field; realigning categories would match calls to where the work
lives. Second, the orphan-organism gap is not self-closing: oil palm, sugarcane,
cassava, and yam remain structurally under-served ($z < -3$, widening since
1990) and need ring-fenced calls and career pathways rather than open
competition with staple-crop and model-organism research. Third, the 2019--2022
sleeping-beauty awakenings on climate-relevant taxa show that latent scientific
capacity can be unlocked quickly when policy and public attention align --- but
that this alignment is fragile and should be institutionalised before the window
closes.
